\BourbakiBackMatter{Bibliography}

{[}1{]} E. ARTIN -- ``Zur Theorie der hypercomplexen Zahlen'', Abh. math. Sem. Univ. Hamburg 5 (1928), p.~251--260; Collected papers, Reading (Addison Wesley), 1965, p.~307--316.

{[}2{]} R. BRAUER -- ``Über Systeme hypercomplexer Zahlen'', Math. Zeitschr. 30 (1929), p.~79--107; Collected papers, vol.~I, Cambridge, Massachusetts (The MIT Press), 1980, p.~40--68.

{[}3{]} W. BURNSIDE - ``On the condition of reducibility of any group of linear substitutions'', Proc. Lond. Math. Soc. 3 (1905), p.~430-434.

{[}4{]} É. CARTAN -- ``Les groupes bilinéaires et les systèmes de nombres complexes'', Ann. Fac. Sc. Toulouse 12 (1898), p.~1--99; Œuvres complètes, partie II, vol.~I, Paris (Gauthier-Villars), 1952, p.~7--105.

{[}5{]} \_\_\_\_ ``Nombres complexes'', Encycl. Sci. Math., éd. française 15 (1908); Œuvres complètes, partie II, vol.~I, Paris (Gauthier-Villars), 1952, p.~107--246.

{[}6{]} A. CAYLEY -- ``On algebraic couples'', Phil. Mag. (3) 27 (1845), p.~38--40; Collected mathematical papers, t. I, Cambridge, 1889, p.~128--131.

{[}7{]} \_\_\_\_ ``On the theory of groups, as depending on the equation \(\theta^{n}=1\) '', Phil. Mag. (4) 7 (1854), p.~40--47; Collected mathematical papers, t. II, Cambridge, 1889, p.~123--130.

{[}8{]} \_\_\_\_``A memoir on the theory of matrices'', Phil. Trans. R. Soc. Lond. 148 (1858), p.~17--37; Collected mathematical papers, t. II, Cambridge, 1889, p.~475--496.

{[}9{]} W. CLIFFORD -- ``Preliminary sketch of biquaternions'', Proc. Lond. Math. Soc. 4 (1873), p.~381--395; Mathematical papers, London (Macmillan), 1882, p.~181--200.

{[}10{]} \_\_\_\_ ``On the classification of geometric algebras'', (1877), Mathematical papers, London (Macmillan), 1882, p.~397--401.

{[}11{]} \_\_\_\_``Applications of Grassmann's extensive algebra'', Amer. Journ. of Math. 1 (1878), p.~350--358; Mathematical papers, London (Macmillan), 1882, p.~266--276.

{[}12{]} R. DEDEKIND -- ``Zur Theorie der aus n Haupteinheiten gebildeten komplexen Grössen'', Gött. Nachr. (1885), p.~141--159; Gesammelte mathematische Werke, t. II, Braunschweig (Vieweg), 1931--32, p.~1--19.

{[}13{]} \_\_\_\_ ``Aus Briefen an Frobenius'', (1886), Gesammelte mathematische Werke, t. II, Braunschweig (Vieweg), 1931--32, p.~414--442.

{[}14{]} \_\_\_\_ ``Über eine Erweiterung des Symbols (a, b) in der Theorie der Moduln'', Gött. Nachr. (1895), p.~183--188; Gesammelte mathematische Werke, t. II, Braunschweig (Vieweg), 1931--32, p.~59--86.

{[}15{]} \_\_\_\_``Über Gruppen, deren sämtliche Teiler Normalteiler sind'', Math. Ann. 48 (1897), p.~548--561; Gesammelte mathematische Werke, t. II, Braunschweig (Vieweg), 1931--32, p.~87--102.

{[}16{]} M. DEURING - Algebren, Erg. der Math., vol.~4, Berlin (Springer), 1935.

{[}17{]} L. DICKSON - ``Linear associative algebras and abelian equations'', Trans. Amer. Math. Soc. 15 (1914), p.~31-46; Mathematical papers, vol.~II, New York (Chelsea), 1975, p.~367-382.

{[}18{]} G. FROBENIUS -- ``Über lineare Substitutionen und bilineare Formen'', Crelle's J. 84 (1878), p.~1--63; Gesammelte Abhandlungen, Band I, Berlin (Springer), 1968, p.~343--405.

{[}19{]} \_\_\_\_ ``Über die Primfactoren der Gruppendeterminante'', Berliner Sitzungsber. (1896), p.~1343--1382; Gesammelte Abhandlungen, Band III, Berlin (Springer), 1968, p.~38--77.

{[}20{]} \_\_\_\_ ``Über Gruppencharaktere'', Berliner Sitzungsber. (1896), p.~985--1021; Gesammelte Abhandlungen, Band III, Berlin (Springer), 1968, p.~1--37.

{[}21{]} \_\_\_\_``Über die Darstellung der endlichen Gruppen durch lineare Substitutionen'', Berliner Sitzungsber. (1897), p.~944--1015; Gesammelte Abhandlungen, Band III, Berlin (Springer), 1968, p.~82--103.

{[}22{]} \_\_\_\_ ``Theorie der hypercomplexen Grössen'', Berliner Sitzungsber. (1903), p.~504--537 and 634--645; Gesammelte Abhandlungen, Band III, Berlin (Springer), 1968, p.~284--329.

{[}23{]} H. GRASSMANN -- ``Sur les différents genres de multiplication'', Crelle's J. 49 (1855), p.~123--141; Gesammelte math. und phys. Werke, t. II, Leipzig (Teubner), 1904, p.~199--217.

{[}24{]} \_\_\_\_ ``Die Ausdehnungslehre von 1862'', (1862), Gesammelte math. und phys. Werke, t. I, Leipzig (Teubner), 1896, p.~1--383.

{[}25{]} W. HAMILTON - Lectures on quaternions, Dublin (Hodges and Smith), 1853.

{[}26{]} O. HÖLDER -- ``Zurückführung einer beliebigen algebraischen Gleichung auf eine Kette von Gleichungen'', Math. Ann. 34 (1889), p.~26--56.

{[}27{]} C. HOPKINS - ``Rings with minimal conditions for left ideals'', Ann. of Math. 40 (1939), p.~712-730.

{[}28{]} W. KRULL -- ``Über verallgemeinerte endliche Abelsche Gruppen'', Math. Zeitschr. 23 (1925), p.~161--196; Gesammelte Abhandlungen, vol.~1, Berlin (de Gruyter), 1999, p.~263--298.

{[}29{]} \_\_\_\_ ``Theorie und Anwendung der verallgemeinerten Abelschen Gruppen'', Sitzungsber. Heidelberger Akad. (1926), p.~1--32; Gesammelte Abhandlungen, vol.~1, Berlin (de Gruyter), 1999, p.~299--328.

{[}30{]} E. LAGUERRE -- ``Sur le calcul des systèmes linéaires'', J. de l'école polytechnique 42 (1867), p.~215--264; Œuvres, t. I, Paris (Gauthier-Villars).

{[}31{]} T. MOLIEN - ``Ueber Systeme höherer complexer Zahlen'', Math. Ann. 41 (1893), p.~83-156.

{[}32{]} \_\_\_\_``Ueber die Invarianten der linearen Substitutionsgruppen'', Berliner Sitzungsber. (1897), p.~1152--1156.

{[}33{]} E. NOETHER -- ``Abstrakter Aufbau der Idealtheorie in algebraischen Zahl- und Funktionenkörpern'', Math. Ann. 96 (1927), p.~26--61; Gesammelte Abhandlungen, Berlin (Springer), 1983, p.~493--528.

{[}34{]} \_\_\_\_ ``Hyperkomplexe Grössen und Darstellungstheorie'', Math. Zeitschr. 30 (1929), p.~641--692; Gesammelte Abhandlungen, Berlin (Springer), 1983, p.~563--614.

{[}35{]} \_\_\_\_ ``Nichtkommutative Algebren'', Math. Zeitschr. 37 (1933), p.~514--541; Gesammelte Abhandlungen, Berlin (Springer), 1983, p.~642--669.

{[}36{]} E. NOETHER et R. BRAUER -- ``Über minimale Zerfällungskörper irreduzibler Darstellungen'', Berliner Sitzungsber. (1927), p.~221--228; Emmy Noether Gesammelte Abhandlungen, Berlin (Springer), 1983, p.~552--559.

{[}37{]} E. NOETHER et W. SCHMEIDLER -- ``Moduln in nichtkommutativen Bereichen'', Math. Zeitschr. 8 (1920), p.~1--35; Emmy Noether Gesammelte Abhandlungen, Berlin (Springer), 1983, p.~318--352.

{[}38{]} B. PEIRCE - ``Linear associative algebra'', Amer. Journ. of Math. 4 (1881), p.~97-221.

{[}39{]} C. PEIRCE - ``On the algebras in which division is unambiguous'', Amer. Journ. of Math. 4 (1881), p.~225-229.

{[}40{]} \_\_\_\_ ``On the relative forms of the algebras'', Amer. Journ. of Math. 4 (1881), p.~221--225.

{[}41{]} H. POINCARÉ -- ``Sur les nombres complexes'', C. R. Acad. Sci. 99 (1884), p.~740--742; Œuvres, t. V, Paris (Gauthier-Villars), 1916--1954, p.~77--79.

{[}42{]} \_\_\_\_ ``Sur l'intégration algébrique des équations linéaires et les périodes des intégrales abéliennes'', Journ. de Math. (5) 9 (1903), p.~139--212; Œuvres, t. III, Paris (Gauthier-Villars), 1916--1954, p.~106--166.

{[}43{]} G. SCHEFFERS -- ``Zurückführung complexer Zahlensysteme auf typische Formen'', Math. Ann. 39 (1891), p.~293--390.

{[}44{]} I. SCHUR -- ``Über eine Klasse von Matrizen, die sich einer gegebenen Matrix zuordnen lassen'', Diss., Berlin, 1901; Gesammelte Abhandlungen, Band I, Berlin (Springer), 1973, p.~1--72.

{[}45{]} \_\_\_\_``Über die Darstellung der endlichen Gruppen durch gebrochene lineare Substitutionen'', Crelle's J. 127 (1904), p.~20--50; Gesammelte Abhandlungen, Band I, Berlin (Springer), 1973, p.~86--116.

{[}46{]} \_\_\_\_ ``Neue Begründung der Theorie der Gruppencharaktere'', Berliner Sitzungsber. (1905), p.~406--432; Gesammelte Abhandlungen, Band I, Berlin (Springer), 1973, p.~143--169.

{[}47{]} \_\_\_\_``Arithmetische Untersuchungen über endliche Gruppen linearer Substitutionen'', Berliner Sitzungsber. (1906), p.~164--184; Gesammelte Abhandlungen, Band I, Berlin (Springer), 1973, p.~177--197.

{[}48{]} J.-A. DE SÉGUIER - Éléments de la théorie des groupes abstraits, Paris (Gauthier-Villars), 1904.

{[}49{]} T. SKOLEM - Zur Theorie des assoziativen Zahlensysteme, Oslo Vid. Akad. Skr. I, 1927.

{[}50{]} J. M. WEDDERBURN - ``A theorem on finite algebras'', Trans. Amer. Math. Soc. 6 (1905), p.~349-352.

{[}51{]} \_\_\_\_``On hypercomplex numbers'', Proc. Lond. Math. Soc. (2) 6 (1908), p.~77--118.

{[}52{]} \_\_\_\_``A type of primitive algebra'', Trans. Amer. Math. Soc. 15 (1914), p.~162--166.

{[}53{]} \_\_\_\_ ``On division algebras'', Trans. Amer. Math. Soc. 22 (1921), p.~129--135.

{[}54{]} K. WEIERSTRASS -- ``Zur Theorie der aus n Haupteinheiten gebildeten complexen Grössen'', Gött. Nachr. (1884), p.~395--414; Mathematische Werke, t. II, Berlin (Mayer und Müller), 1895, p.~311--332.

\BourbakiBackMatter{Notation Index}

\(A_{M}\) \ldots.. 1 K{[}G{]} \ldots.. 1 A{[}X{]} \ldots.. 9 A{[}X{]} \(_{\sigma,d}\) \ldots.. 11 {[}M:L{]} \ldots.. 32 IsA\{X,Y\} \(_{\ddagger}\) \ldots.. 51 cl(X) \ldots.. 51 F(A) \ldots.. 51 S(A) \ldots.. 51 s \(^{0}\) , d \(^{0}\) \ldots.. 53 T(V) \ldots.. 58 H(M) \ldots.. 58 αM \ldots.. 59 βV \ldots.. 59 MS \ldots.. 65 αM \ldots.. 69 τ(M) \ldots.. 80 S\_M \ldots.. 84 {[}B:A{]} \(_{s}\) \ldots.. 124 R\_A(M) \ldots.. 151 R(M) \ldots.. 151 R(A) \ldots.. 154 K(C) \ldots.. 186 {[}E{]}C,{[}E{]} \ldots.. 186 K'(C) \ldots.. 188 {[}E{]}'C \ldots.. 189 \(\mathrm{K}^{\prime}(\mathcal{C})^{+}\) 189 \(\ell_{\mathrm{S}}(\mathrm{E})\) 189 \(\mathcal{L}\mathcal{F}(\mathrm{A})\) 191 \(\mathrm{R(A)}\) 191 \(\mathrm{R_K(A)}\) 194 \(\mathrm{R_K(G)}\) 198 \(\mathcal{R}(\mathfrak{g})\) 199 \(\mathrm{K_0(A)}\) 199 \(i(f)\) 204 \(i(\mathrm{B,A})\) 204 \(h(f)\) 204 \(h(\mathrm{B,A})\) 204 \(\mathrm{C^n(A,P)}\) 239 \(\mathrm{H^n(A,P)}\) 240 \(\mathrm{A^f}\) 254 \(\mathrm{Iso}_{\mathrm{K}}\{\mathrm{A},\mathrm{B}\}\) 277 \(\operatorname {cl}(A)\) 277 \(\operatorname {cl}_{\mathrm{K}}(A)\) 277\\
Br(K) 278 \(\operatorname {Ex}_{\tau}(\mathrm{G},\mathrm{F})\) 286 \(\mathcal{I}_{\tau}\) 286 \(\mathrm{Z}^2 (\mathrm{G},\mathrm{F})\) 295 \(\mathrm{C}^1 (\mathrm{G},\mathrm{F})\) 295 \(\mathrm{H}^2 (\mathrm{G},\mathrm{F})\) 296 \(\operatorname {Res}_{\mathrm{G}'}^{\mathrm{G}}\) 299 \(\operatorname {Cor}_{\mathrm{H}}^{\mathrm{G}}\) 303

\(\mathrm{Coind}_{\mathrm{H}}^{\mathrm{G}}(\mathrm{E})\) 305 \(\mathbf{A}[\mathcal{E};\mathrm{L}]\) 315 \(\mathrm{H}^n (\mathrm{G},\mathrm{M})\) 324 \(\mathrm{Res}^q\) 325 \(\mathrm{Inf}^q\) 325 \(\mathrm{Cor}^q\) 326 \(\mathrm{Br(A)}\) 330 \(\mathrm{Br(B / A)}\) 330 \(\mathrm{Pcrd}_{\mathrm{A / K}}(a;\mathrm{X})\) 340 \(\mathrm{Trd}_{\mathrm{A / K}}\) 340 \(\mathrm{Nrd}_{\mathrm{A / K}}\) 340 \(\mathrm{Pm}_{\mathrm{A / K}}(a;\mathrm{X})\) 351 \(\Theta (\mathrm{A})\) 376 \(c_{\mathrm{E}}(x,x^{*})\) 377 \(\Theta^{\mathrm{ss}}(\mathrm{A})\) 380 \(\mathrm{Hom}_{\mathrm{G}}(\pi ,\pi ')\) 398

\(\mathcal{Z}_{\mathrm{K}}(\mathrm{G})\) 398 \(\chi_{\pi}\) 398 \(\pi \otimes \pi^{\prime}\) 399 \(\pi \boxtimes \pi^{\prime}\) 400 \(\operatorname{Res}_{\mathrm{H}}^{\mathrm{G}}(\pi)\) 402 \(\operatorname{Ind}_{\mathrm{H}}^{\mathrm{G}}(\sigma)\) 402 \(\operatorname{Coind}_{\mathrm{H}}^{\mathrm{G}}(\sigma)\) 403 \(\widehat{\mathbf{G}}\) 406 \(\mathrm{F}(\widehat{\mathrm{G}})\) 407 \(\det u\) 452 \(\det(A)\) 453 \(\mathbf{SL}_n(\mathrm{D})\) 455 \(\mathbf{SL}(\mathrm{V})\) 457 \(\operatorname{sdet} X\) 460 \(\operatorname{Hom}_{\mathrm{A}}^{\mathrm{f}}(\mathrm{E},\mathrm{F})\) 463 \(\operatorname{End}_{\mathrm{A}}^{\mathrm{f}}(\mathrm{E})\) 463

\BourbakiBackMatter{Terminology Index}

2-coboundary, 295\\
2-cocycle, 295

\BourbakiIndexLetter{A}

Absolutely semisimple algebra, 231 module, 229 Absolutely without radical, module, 246 Adapted basis, 460 Additive function of modules, 183 set of classes of modules, 183 weakly --- function of modules, 183 Adjunction isomorphism, 59 Adverse, 441 Algebra absolutely semisimple, 231 algebraic, 359 Artinian, 444 Azumaya, 271 central, 251 coinduced, 305 Galois, 310 (K, G)- ---, 304 primitive, 443 quaternion, 361 semisimple, 136 similar, 280 simple, 120

Artinian\\
algebra, 444\\
module, 1\\
ring, 4\\
Autoinjective\\
ring, 53\\
Azumaya algebra, 271\\
Azumaya's theorem, 32

\BourbakiIndexLetter{B}

Balanced module, 77\\
Basis, adapted, 460\\
Bicommutant, 77\\
Bimodule\\
invertible, 100\\
left, 56\\
Blocks, 180\\
Brauer group, 279, 330\\
Burnside's theorem, 84, 429

\BourbakiIndexLetter{C}

Calculus of fractions, 17\\
Cartan matrix, 201\\
Cayley (--- - Hamilton theorem), 338\\
Central\\
algebra, 251\\
function, 398\\
Centralizer, 77\\
Character

Brauer, 434\\
of a representation, 382, 398\\
Chevalley (--- --Warning theorem), 355\\
Class of a module, 51\\
Coboundary, 295\\
Cocycle, 295\\
associated with a \(\tau\) -extension, 296\\
Coefficient\\
of a module, 377\\
of a representation, 377\\
Cogebra, sub---, 391\\
Cohomology\\
group, 324\\
of algebras, 239\\
Coinduced\\
algebra, 305\\
representation, 403\\
Colength of a module, 53\\
Commutant, 77\\
Compatible homomorphisms, 324\\
Component, isotypical --- of type S, 65\\
Corestriction, 303, 326\\
Cotranspose, 338\\
Countermodule, 8\\
Cover, projective, 161\\
Criterion, Mackey's, 426\\
Cross product, 315

\BourbakiIndexLetter{D}

Decomposable idempotent, 146\\
Decomposition, Weyr-Fitting, 27\\
Definition of fractions, 18\\
Degenerate, non---subalgebra, 370\\
Degree left, 124\\
of a representation, 373, 397\\
over a subring, 124, 204\\
reduced --- of a simple algebra, 253\\
Dense subring, 129\\
Density theorem, 88\\
Density, Jacobson's --- theorem, 442\\
Derivation, 9\\
inner, 240\\
Description canonical, 62, 70

of a module, 69\\
of an isotypical module, 62\\
Determinant, 452, 453\\
Diagram, Young, 430\\
Direct image of a \(\tau\) -extension, 289\\
Divisible, cardinal --- by, 123

\BourbakiIndexLetter{E}

Element idempotent, 145 integral, 428 p-regular, 433 without multiple factors, 155 Equivalence, Morita --- of algebras, 100 of modules, 103 Equivalent idempotents, 40 Essential submodule, 75 Exponent, 322 Extension of an algebra, 248 product of, 293 splitting field, 281 τ-extension, 285 product of, 293 semitrivial, 286

\BourbakiIndexLetter{F}

Factor, simple --- of type S, 141 Faithful representation of an algebra, 374 Field C1, 358 finite, 357 splitting, 281 Field, residue, 26 Fitting, (Weyr--- decomposition), 27 Flat, absolutely --- ring, 171 Foot, 130 Form, trace, 464 Formula, Frobenius, 432 Fourier inversion formula, 408 Fractions calculus of, 17 definition of, 18 monoid of, 17

ring of, 18\\
Frobenius\\
formula, 432\\
reciprocity, 203, 403\\
ring, 53\\
Function\\
additive, of modules, 183\\
central, 398\\
weakly additive, of modules, 183

\BourbakiIndexLetter{G}

Galois algebra, 310 group of automorphisms, 274 subring, 274 Gelfand (--- --Mazur theorem), 368 Generating module, 79 Goldie's theorem, 149 Grothendieck group, 186, 189 ring, 198 Group Brauer, 279, 330 Galois, of automorphisms, 274 Grothendieck, 186 for direct sums, 189 unimodular, 455, 457

\BourbakiIndexLetter{H}

Hamilton (Cayley- --- theorem), 338\\
Height\\
of a morphism, 204\\
of a subring, 204\\
Hereditary set of classes of modules, 183\\
Hilbert's\\
90 theorem, 327\\
Nullstellensatz, 461\\
Homomorphism\\
compatible, 324\\
corestriction, 303, 326\\
inflation, 299, 325\\
linking, 324\\
of finite rank, 463\\
restriction, 299, 325\\
Homothety ring, 1

\BourbakiIndexLetter{I}

Ideal maximal, 26, 141, 435 minimal, 50 nilpotent, 149, 156 prime, 133 primitive, 133 regular, 435 semiprime two-sided, 149 trace, 80

Idempotent central, of a group algebra, 408 decomposable, 146 indecomposable, 146

Idempotents, orthogonal, 146

Indecomposable idempotent, 146 module, 30

Index, 322 of a morphism, 204 of a subring, 204

Induced representation, 402

Inflation homomorphism, 299, 325

Intertwining operator, 398

Inverse image of a cocycle, 299 of a \(\tau\) -extension, 287

Inverse of a module, 100

Invertible bimodule, 100

Isomorphism, adjunction, 59

Isotypical component of type S, 65 module, 61

\BourbakiIndexLetter{J}

Jacobson density theorem, 88 radical, 154 Jacobson's theorem, 156

\BourbakiIndexLetter{K}

Krull--Remak--Schmidt theorem, 37

\BourbakiIndexLetter{L}

Lemma

Nakayama's, 158\\
Schur's, 47\\
Length\\
of a module, 72\\
of a ring, 7\\
Linked modules, 180\\
Linking homomorphism, 324\\
Local endomorphism ring, module with a, 31\\
Local ring, 26

\BourbakiIndexLetter{M}

Mackey's criterion, 426\\
Maschke's theorem, 401\\
Matrix, Cartan, 201\\
Maximal\\
ideal, 26, 141, 435\\
submodule, 48\\
Mazur (Gelfand- --- theorem), 368\\
Mean, 325\\
Minimal, ideal, 50\\
Module\\
absolutely semisimple, 229\\
absolutely without radical, 246\\
Artinian, 1\\
balanced, 77\\
generating, 79\\
indecomposable, 30\\
isotypical, 61\\
Noetherian, 1\\
of a representation, 373\\
of finite colength, 53\\
of type \(\mathcal{C}\) , 183\\
primordial, 31\\
semi-Artinian, 74\\
semiprimordial, 32\\
semisimple, 55\\
simple, 45\\
with local endomorphism ring, 31\\
Modules\\
linked, 180\\
stably isomorphic, 200\\
Monoid of fractions, 17\\
Morita equivalence\\
of algebras, 100

of modules, 103\\
Multiplicity, 72, 189\\
primordial, 34\\
Multiplier\\
of a pseudobimodule, 445\\
of an algebra, 445

\BourbakiIndexLetter{N}

Nakayama's lemma, 158\\
Nil ideal, 156\\
Nilpotent ideal, 156\\
Nilradical, 157\\
Noether (Skolem- --- theorem), 256\\
Noetherian module, 1 ring, 4\\
Norm, reduced, 340\\
n-regular ring, 178

\BourbakiIndexLetter{O}

Operator, intertwining, 398 Orthogonal idempotents, 146 of an ideal, 53 Orthogonality relation for characters, 410 Schur --- relation, 410 second --- relation for characters, 412

\BourbakiIndexLetter{P}

Partition of an idempotent, 146\\
Permutation representation, 425\\
Polynomial\\
principal, 351\\
reduced characteristic, 340\\
Prime\\
ideal, 133\\
radical, 167\\
Primitive\\
ideal, 133\\
ring, 120\\
Primitive algebra, 443\\
Primordial module, 31\\
Principal polynomial, 351\\
Product

cross, 315\\
external tensor, 400\\
tensor, of representations, 399\\
Projective cover, 161\\
Pseudobimodule, 445\\
Pseudomodule\\
left, right, 439\\
semisimple, 442\\
simple, 440\\
Pseudoregular ring, 178

\BourbakiIndexLetter{Q}

Quasi-simple ring, 120\\
Quaternions, 361

\BourbakiIndexLetter{R}

Radical Jacobson, 154 of a module, 151 of a ring, 154 of an algebra, 440 prime, 167 Reciprocity, Frobenius, 203, 403 Reduced, degree, 253 Regular ideal, 435 ring, 171 Related representations, 416 Relation orthogonality, for characters, 410 Schur orthogonality, 410 second orthogonality, for characters, 412 Representation biregular, 399 coinduced, 403 completely reducible, 374 contragredient, 399 dual, 399 faithful, 374 induced, 402 irreducible, 374 linear of a group, 397 of an algebra, 373

matrix, 377\\
monomial, 427\\
permutation, 425\\
quotient, 374\\
regular, 375, 399\\
semisimple, 374\\
simple\\
of a group, 398\\
of an algebra, 374\\
transposed, 374\\
unit, 399\\
unitary, 421\\
Restricted dual, 376\\
Restriction homomorphism, 299, 325\\
Ring\\
absolutely flat, 171\\
Artinian, 4\\
autoinjective, 53\\
Frobenius, 53\\
Grothendieck, 198\\
homothety, 1\\
local, 26\\
Noetherian, 4 \(n\) -regular, 178\\
of fractions, 18\\
polynomial, 11\\
primitive, 120\\
pseudoregular, 178\\
quasi-simple, 120\\
semiprime, 149\\
semisimple, 136\\
simple, 120\\
von Neumann regular, 171\\
without radical, 154\\
Zorn, 171

\BourbakiIndexLetter{S}

Saturation, left, right, 19\\
Schur orthogonality relation, 410\\
Schur's lemma, 47\\
Second orthogonality relation for characters, 412\\
Semi-Artinian module, 74\\
Semiprime ring, 149\\
Semiprimordial module, 32

Semisimple algebra, 136 module, 55 pseudomodule, 442 ring, 136

Semisimplification, 191 Semitrivial τ-extension, 286 Set of classes of modules additive, 183 hereditary, 183 Similar algebras, 280 Simple algebra, 120 factor of type S, 141 pseudomodule, 440 representation of a group, 398 ring, 120 Simple module, 45 Skolem (--- --Noether theorem), 256 Socle, 65, 130 right, left, 130 Splitting field, 281 Subalgebra commutative maximal, 261 maximal commutative semisimple, 264 maximal étale, 264 nondegenerate, 370 Subcogebra, 391 Submodule essential, 75 maximal, 48 Sub-pseudomodule, 439 Subrepresentation, 374 Subring dense, 129 Galois, 274 weakly Galois, 274 Support of a semisimple module, 66, 190

\BourbakiIndexLetter{T}

Theorem

Azumaya's, 32\\
Burnside's, 84, 429\\
Cayley--Hamilton, 338\\
Chevalley--Warning, 355\\
Gelfand--Mazur, 368\\
Goldie, 149\\
Hilbert's 90, 327\\
Hilbert's Nullstellensatz, 461\\
Jacobson's, 156\\
Jacobson's density, 88, 442\\
Krull--Remak--Schmidt, 37\\
Maschke's, 401\\
Skolem--Noether, 256\\
Tsen's, 358, 360\\
Wedderburn's, 120, 135, 245, 357\\
Trace\\
form, 464\\
ideal, 80\\
of a representation, 398\\
reduced, 340\\
Tracial mapping, 115\\
Transvection, 457\\
Tsen's theorem, 358, 360

\BourbakiIndexLetter{U}

Unimodular group, 455\\
Unit, left, right, 435

\BourbakiIndexLetter{W}

Warning (Chevalley--- theorem), 355 Weakly Galois subring, 274 Wedderburn's theorem, 120, 135, 245, 357 Weyr (--- -Fitting decomposition), 27

\BourbakiIndexLetter{Y}

Young diagram, 430

\BourbakiIndexLetter{Z}

Zorn ring, 171
